\documentclass[answers,12pt,addpoints]{exam} 
\usepackage{import}

\import{C:/Users/prana/OneDrive/Desktop/MathNotes}{style.tex}

% Header
\newcommand{\name}{Pranav Tikkawar}
\newcommand{\course}{Spatial Stats}
\newcommand{\assignment}{Homework 7.1}
\author{\name}
\title{\course \ - \assignment}

\begin{document}
\maketitle

\begin{questions}
\question Microergodicity is a very helpful property! I tried looking into the reference (117) but I didnt know where to look. I found this paper https://arxiv.org/html/2603.23959v1 but I was struggling to parse it. I will read it later and potentially make comments. 

\question I understasnd better no why the statment in class that I (which built off Liron's) about how we can treat this test with some $\alpha$ level and see if our $R_m$ reach a level where the overlap of the probability laws is only $\alpha$ is not super useful, but somewhat interesting. There is a stronger condition that needs to be met to get a stronger conclusion. 

\question For the periodic case of GPs the equivalence is governed by the smoothness/rate of decay of the fourier coefficients. This makes sense when we consider the image we viewed in class since fast rates of decay in frequency are slow in space, and we require the measures in their limit to not have any support on each other. 

\question I am struggling a bit with chapter 8.5. I think it is mainly about extending the periodic theory to a theory with spectral denisty. It seems to struggle on finite bounded data. I will need to re-read to better understand it. 

\question For dim $d \leq 3$ the set of parameters $\phi, \kappa, A$ are microergodic and consistently estimable with fixed domain asymptotes. It is mentioned that there is no rigourous proof that a sequence of MLE estimates of $\phi,\kappa$ are consistent. Does this mean there some classes of data or some specific methods that work, or historically has it show to be true by experiment. 

\question I will read these parts again tomorrow morning and review your feedback, since I feel less confident in my understanding of these parts. 


\end{questions}

\end{document}